% ============================================================
% donghx 格点 QCD 算法实现 — 物理与代码联合报告
% 编译：xelatex 两遍；16:9 横板；主体算法采用单列 tabular 行式表达
% 本报告只使用当前目录代码与已有文档的可追溯证据，不填入未运行的数值结果。
% ============================================================
\documentclass[UTF8,aspectratio=169,11pt]{ctexbeamer}

\usepackage{amsmath,amssymb}
\usepackage{graphicx}
\usepackage{booktabs}
\usepackage{tabularx}
\usepackage{array}
\usepackage{tikz}
\usetikzlibrary{arrows.meta,positioning,calc,fit,shapes.geometric}
\usepackage{listings}
\usepackage{xcolor}
\usepackage{microtype}
\usepackage{hyperref}

\definecolor{primary}{HTML}{17365D}
\definecolor{accent}{HTML}{1F77B4}
\definecolor{evidence}{HTML}{1E8449}
\definecolor{warning}{HTML}{C55A11}
\definecolor{muted}{HTML}{5B6573}
\definecolor{panel}{HTML}{F4F7FA}
\definecolor{linegray}{HTML}{CBD5E1}
\definecolor{good}{HTML}{EAF5EE}
\definecolor{warnbg}{HTML}{FFF4E8}

\setbeamersize{text margin left=0.55cm,text margin right=0.55cm}
\setlength{\emergencystretch}{2em}
\setbeamertemplate{navigation symbols}{}
\setbeamertemplate{blocks}[rounded][shadow=false]
\setbeamercolor{normal text}{fg=black!86,bg=white}
\setbeamercolor{structure}{fg=primary}
\setbeamercolor{frametitle}{fg=primary,bg=white}
\setbeamercolor{block title}{fg=primary,bg=panel}
\setbeamercolor{block body}{fg=black!86,bg=panel}
\setbeamerfont{frametitle}{size=\large,series=\bfseries}
\setbeamerfont{normal text}{size=\normalsize}
\setbeamertemplate{frametitle}{%
  \vskip0.12cm\insertframetitle\par\vskip0.08cm}
\setbeamertemplate{footline}{%
  \hbox{\begin{beamercolorbox}[wd=\paperwidth,ht=2.3ex,dp=0.9ex,
    leftskip=0.55cm,rightskip=0.55cm]{author in head/foot}%
    \scriptsize\color{muted}\insertshortauthor\hfill
    \insertshorttitle\hfill\insertframenumber/\inserttotalframenumber
  \end{beamercolorbox}}}

\newcommand{\source}[1]{%
  \vspace{0.04em}\par{\raggedright\scriptsize\color{muted}来源：#1\par}}
\newcommand{\verified}{\textcolor{evidence}{\textbf{确证}}}
\newcommand{\inferred}{\textcolor{accent}{\textbf{推断}}}
\newcommand{\unverified}{\textcolor{warning}{\textbf{未验证}}}
\newcommand{\code}[1]{\texttt{\detokenize{#1}}}
\newcommand{\rowlabel}[1]{\textcolor{primary}{\textbf{#1}}}

\lstset{
  basicstyle=\ttfamily\scriptsize,
  backgroundcolor=\color{panel},frame=single,
  breaklines=true,breakatwhitespace=false,keepspaces=true,
  columns=flexible,firstnumber=auto,
  keywordstyle=\color{primary}\bfseries,
  commentstyle=\color{evidence!70!black},
  stringstyle=\color{accent},
  numbers=left,numberstyle=\tiny\color{muted}
}

\title[donghx 格点 QCD 算法]{donghx 格点 QCD 算法实现}
\subtitle{从非连通胶子 OPE 到动量涂抹蒸馏质子 2pt}
\author[代码—物理联合审阅]{代码—物理联合审阅}
\institute{PyQCD / refer/donghx}
\date{2026 年 8 月 28 日}

\begin{document}

% S01
\begin{frame}[plain]
  \titlepage
\end{frame}

% S02
\begin{frame}[t]{本阶段的核心判断：代码已形成两条可复用的底层计算链}
  \begin{center}
    \small
    \begin{tabularx}{0.96\textwidth}{@{}>{\raggedright\arraybackslash}p{0.18\linewidth}
      >{\raggedright\arraybackslash}p{0.34\linewidth}
      >{\raggedright\arraybackslash}X@{}}
      \toprule
      状态 & 已确认的实现 & 对物理目标的意义 \\
      \midrule
      \verified\ 算法骨架
        & 规范场读入、Clover 场强、Wilson 线、颜色迹/空间求和
        & 可计算非定域胶子算符 $\mathcal O(z)$，并保留 $\pm z$ 变体。 \\
      \verified\ 2pt 链
        & 本征矢量 + 动量相位 + perambulator + VVV + 两项 Wick 缩并
        & 为高动量核子态提供 $C_2(t,\vec P)$，作为断连三点链的核子因子。 \\
      \verified\ 工程策略
        & CuPy、\code{opt_einsum}、按 $t_{\rm source}$ 读入、$x$ 向分块
        & 把 $N_{\rm ev}^3$ 中间张量的显存压力转化为可管理的循环。 \\
      \unverified\ 数值结论
        & 当前工作区没有外部 ILDG、eigenvector、perambulator 数据
        & 本报告验证代码结构与语义；不宣称质量、矩阵元或 PDF 数值已复现。 \\
      \bottomrule
    \end{tabularx}
  \end{center}
  \vspace{0.08cm}
  \begin{block}{需要保持的物理边界}
    “$C_3\simeq C_2\times\mathrm{OPE}$” 是代码组织上的两路估计量分离；
    真正的断连矩阵元仍需在组态集合上做相关平均、真空扣除和归一化。
  \end{block}
  \source{\code{docs/AGENTS.md}；\code{docs/donghx_code_analysis.tex:119--189}；
  当前目录输入数据缺失为本次检查实测。}
\end{frame}

% S03
\begin{frame}[t]{物理目标：格点非定域矩阵元通过 LaMET 连接胶子 PDF}
  \begin{columns}[T,totalwidth=\textwidth]
    \begin{column}{0.55\textwidth}
      \begin{block}{物理定义}
        \begin{align}
          \widetilde h(z,P_z)
          &= \frac{\langle P|\mathcal O_{\mu\nu;\rho\sigma}(z)|P\rangle}
                  {\langle P|P\rangle}, \\
          \mathcal O_{\mu\nu;\rho\sigma}(z)
          &= F_{\mu\nu}(0)\,U(0,z\hat n)\,
             F_{\rho\sigma}(z\hat n)\,U(z\hat n,0), \\
          \widetilde g(x,P_z)
          &= \int_{-\infty}^{\infty}\frac{dz}{2\pi}
             e^{ixP_z z}\widetilde h(z,P_z).
        \end{align}
      \end{block}
      \begin{block}{极限与量纲检查}
        $z=0$ 时 Wilson 线退化为单位元；$a\to0$ 时 Clover 场强应趋近连续
        $F_{\mu\nu}$；$P_z$ 越大，LaMET 的幂修正应按
        $\Lambda_{\rm QCD}^2/P_z^2$ 抑制。
      \end{block}
    \end{column}
    \begin{column}{0.41\textwidth}
      \begin{tabular}{@{}l@{}}
        \rowlabel{代码对象与物理对象} \\[0.08cm]
        \code{gauge[t,z,y,x,mu,a,b]} \\
        \quad $\longleftrightarrow U_\mu(x)$ \\[0.05cm]
        \code{pla[mu,nu,t,z,y,x,a,b]} \\
        \quad $\longleftrightarrow F_{\mu\nu}(x)$ \\[0.05cm]
        \code{operators_new_z0_mu2(...)} \\
        \quad $\longleftrightarrow F\,U^\dagger\,\widetilde F\,U$ \\[0.05cm]
        \code{ops[t,dz]} \\
        \quad $\longleftrightarrow$ 时间片上的空间和 \\[0.05cm]
        \code{twopt_slice_pp[ts,ti]} \\
        \quad $\longleftrightarrow C_2^{P^+}(t_s,t_i)$
      \end{tabular}
      \source{\code{docs/donghx_code_analysis.tex:119--189}；
      \code{Operator.py:48--184,339--368}。}
    \end{column}
  \end{columns}
\end{frame}

% S04
\begin{frame}[t]{总流程：同一规范组态分叉为 OPE 分支与核子 2pt 分支}
  \begin{center}
    \small
    \begin{tabularx}{0.95\textwidth}{@{}>{\raggedright\arraybackslash}p{0.19\linewidth}
      >{\centering\arraybackslash}p{0.06\linewidth}
      >{\raggedright\arraybackslash}X@{}}
      \toprule
      阶段 & & 计算状态与物理接口 \\
      \midrule
      输入
        & $\Downarrow$
        & ILDG 原始规范场 $\to$ 复数张量
          \code{[Nt,Nx,Nx,Nx,4,3,3]}。 \\
      OPE 分支
        & $\Longrightarrow$
        & 四叶草 $P_{\mu\nu}\to F_{\mu\nu}$；对偶张量
          $\widetilde F_{\mu\nu}=\frac12\epsilon_{\mu\nu\rho\sigma}F_{\rho\sigma}$；
          Wilson 线连接 $\pm z$ 两端。 \\
      2pt 分支
        & $\Longrightarrow$
        & 低模 $v^{(k)}$ 乘动量相位；读取 $\tau$ 与 VVV；
          颜色、Dirac、本征模指标做费米子反对称缩并。 \\
      断连接口
        & $\Downarrow$
        & 组态级 $O_{\rm gluon}(z,t)$ 与 $C_2(t,\vec P)$
          进入同一统计分析；必须保留协方差与真空扣除。 \\
      下游物理
        & $\Downarrow$
        & $\widetilde h(z,P_z)\to$ Fourier/quasi-PDF $\to$ 匹配；
          该目录只直接覆盖前端算符与核子 2pt。 \\
      \bottomrule
    \end{tabularx}
  \end{center}
  \source{流程为代码依赖关系示意；\code{Calc_ope_unpol.py:94--136}；
  \code{2pt_proton_Cg5gmu_L32x64_mom2_xdir_gpu.py:172--236,349--447}。}
\end{frame}

% S05
\begin{frame}[t]{Clover 场强：四个定向 plaquette 把局部规范几何变成 $F_{\mu\nu}$}
  \small
  \begin{columns}[T,totalwidth=\textwidth]
    \begin{column}{0.54\textwidth}
      \begin{block}{连续定义与离散实现}
        \begin{align}
          P_{\mu\nu}(x)
          &=U_\mu(x)U_\nu(x+\hat\mu)
            U_\mu^\dagger(x+\hat\nu)U_\nu^\dagger(x),\\
          Q_{\mu\nu}
          &=P_{\mu\nu}+P_{\nu,-\mu}
            +P_{-\mu,-\nu}+P_{-\nu,\mu},\\
          F_{\mu\nu}^{\rm clover}
          &\propto -i\,[Q_{\mu\nu}-Q_{\mu\nu}^\dagger].
        \end{align}
        四叶草的中心对称性抵消一阶方向误差；代码以数组滚动实现周期边界。
      \end{block}
    \end{column}
    \begin{column}{0.42\textwidth}
      \begin{tabular}{@{}l@{}}
        \rowlabel{Clover 算法（实现语义）} \\[0.08cm]
        $\text{Input: }U[t,z,y,x,\mu,a,b]$ \\[0.03cm]
        $\text{Place four shifted views of }U$ \\[0.03cm]
        \quad $P_{\mu\nu},P_{\nu,-\mu},P_{-\mu,-\nu},
        P_{-\nu,\mu}$ \\[0.03cm]
        $\text{Build }Q_{\mu\nu}-Q_{\mu\nu}^{\dagger}$ \\[0.03cm]
        $\text{Return }-i(Q-Q^\dagger)/8$ \\[0.03cm]
        $\text{for }\mu,\nu=0,\ldots,3\text{; set }\mu=\nu\text{ to zero}$ \\[0.03cm]
        $\widetilde F_{\mu\nu}
        =\frac12\epsilon_{\mu\nu\rho\sigma}F_{\rho\sigma}$ \\[0.03cm]
        $\text{Output: }[4,4,N_t,N_x,N_x,N_x,3,3]$
      \end{tabular}
    \end{column}
  \end{columns}
  \vspace{0.10cm}
  \begin{block}{实现边界}
    \verified\ 四个方向的反厄米组合与 $\epsilon_{\mu\nu\rho\sigma}$ 张量有代码证据；
    \unverified\ 连续公式中的显式无迹投影在 \code{Operator.py:150--160} 未单独出现，
    因而不能仅凭函数名宣称代码已完成 $SU(3)$ traceless 投影。
  \end{block}
  \source{\code{Operator.py:77--184}；\code{Calc_pla.py:29--52}；
  物理公式参照 \code{docs/donghx_code_analysis.tex:271--485}。}
\end{frame}

% S06
\begin{frame}[t]{非定域 OPE：两条反向 Wilson 线把两端场强闭合成规范不变量}
  \small
  \begin{columns}[T,totalwidth=\textwidth]
    \begin{column}{0.47\textwidth}
      \begin{tabular}{@{}l@{}}
        \rowlabel{算符构造伪代码（单列行式）} \\[0.08cm]
        $\text{Input: }U,\;P_{\mu\nu},\;\widetilde P_{\mu_2\nu_2},
        \delta z,\hat n$ \\[0.04cm]
        $\text{Initialize: }A\gets\mathrm{roll}(P_{\mu\nu},-\delta z,\hat n)$ \\[0.04cm]
        \quad $k=0,\ldots,\delta z-1:
        \;A\gets A\,U_{\hat n}^{\dagger}(\delta z-1-k)$ \\[0.04cm]
        $\text{Insert: }A\gets A\,\widetilde P_{\mu_2\nu_2}(0)$ \\[0.04cm]
        \quad $\text{for }k=0,\ldots,\delta z-1:
        \;A\gets A\,U_{\hat n}(k)$ \\[0.04cm]
        $\text{Trace: }s(t,\vec x)\gets\mathrm{Tr}_{c}A$ \\[0.04cm]
        $\text{Sum: }O(t,\delta z)\gets\sum_{\vec x}s(t,\vec x)$ \\[0.04cm]
        $\text{Output: }O\in\mathbb C^{N_t}$
      \end{tabular}
      \source{伪代码对应 \code{Operator.py:339--368}；
      $\hat n$ 由 \code{zdir} 选择。}
    \end{column}
    \begin{column}{0.49\textwidth}
      \begin{block}{物理对应}
        \[
          \begin{aligned}
          O_{\mu\nu;\mu_2\nu_2}(z)
          &=\sum_{\vec x}\mathrm{Tr}_{c}\bigl[
          U^\dagger F_{\mu\nu}(z)U\widetilde F_{\mu_2\nu_2}(0)\bigr].
          \end{aligned}
        \]
        代码中的 \code{np.roll} 是格点平移，\code{einsum} 是颜色矩阵乘法；
        最后的颜色迹和横向空间求和把矩阵对象变成每个时间片的标量。
      \end{block}
      \begin{block}{不变量检查}
        $\delta z=0$：两段 Wilson 线循环为空，退化为局部场强乘积；
        $U^\dagger U=1$：规范变换因子在颜色迹中闭合；
        方向翻转：负向函数改变平移与链路方向。
      \end{block}
    \end{column}
  \end{columns}
  \source{\code{Operator.py:339--399}。}
\end{frame}

% S07
\begin{frame}[t]{Lorentz 分量与 MPI：并行切分的是物理指标，不是同一算符的重复计算}
  \small
  \begin{center}
    \begin{tabularx}{0.95\textwidth}{@{}>{\raggedright\arraybackslash}p{0.14\linewidth}
      >{\centering\arraybackslash}p{0.17\linewidth}
      >{\raggedright\arraybackslash}X
      >{\raggedright\arraybackslash}p{0.27\linewidth}@{}}
      \toprule
      MPI rank & 指标 $(\mu,\nu)$ & $z\parallel\hat n$ 时的物理分量 & 输出/用途 \\
      \midrule
      0 & $(3,z+1)$ & $F_{t x}$ 或 $F_{t y}$：电场型分量
        & unpol/helicity 的第一路 \\
      1 & $(3,z+2)$ & 另一电场型分量
        & 与 rank 0 形成横向对称配对 \\
      2 & $(z+1,z+2)$ & $F_{xy}$：空间磁场型分量
        & 补齐第三个独立分量 \\
      \midrule
      \multicolumn{4}{l}{\code{Calc_ope_unpol.py}：传入 $(P,P)$，得到 $FUFU^\dagger$ 型算符。}\\
      \multicolumn{4}{l}{\code{Calc_ope_helicity.py}：传入 $(P,\widetilde P)$，同时保存 $+z$ 与 $-z$。}\\
      \bottomrule
    \end{tabularx}
  \end{center}
  \vspace{0.12cm}
  \begin{tabular}{@{}l@{}}
    $\text{for each rank }r\in\{0,1,2\}:\quad
      (\mu,\nu,\mu_2,\nu_2)\gets\text{rank-specific assignment}$ \\
    $\quad \text{read }U\;\longrightarrow\;P\;(\longrightarrow\widetilde P)
      \;\longrightarrow\;\{O(t,\delta z)\}_{\delta z=0}^{\delta z_{\max}-1}
      \;\longrightarrow\;\texttt{npy}$ \\
    $\text{Boundary branch: }\texttt{link\_dir}\in\{\texttt{x,y,z}\}
      \;\Longrightarrow\; \hat n\in\{\hat x,\hat y,\hat z\}$
  \end{tabular}
  \vspace{0.08cm}
  \unverified\ 三分量的下游归一化、统计权重与最终物理投影不在本目录中。
  \source{\code{Calc_ope_unpol.py:53--136}；
  \code{Calc_ope_helicity.py:50--142}。}
\end{frame}

% S08
\begin{frame}[t]{Distillation + momentum smearing：低模子空间承载可控的高动量核子态}
  \small
  \begin{columns}[T,totalwidth=\textwidth]
    \begin{column}{0.51\textwidth}
      \begin{block}{物理展开}
        \[
          \begin{aligned}
          S(x,y)&\simeq\sum_{k,l=1}^{N_{\rm ev}}
          v^{(k)}(x)\,\tau^{(k,l)}(t_x,t_y)\\
          &\quad\times v^{(l)\dagger}(y),\\
          v_{\rm mom}^{(k)}(x)&=e^{-i\boldsymbol\zeta\cdot\boldsymbol x}v^{(k)}(x).
          \end{aligned}
        \]
        $N_{\rm ev}$ 控制空间 Laplacian 低模截断；相位把源/汇波函数的 overlap
        推向目标动量，而不直接改变规范场。
      \end{block}
      \begin{block}{量纲与极限}
        $\boldsymbol\zeta=\boldsymbol 0$ 恢复普通 distillation；
        相位指数无量纲，代码用 $2\pi/N_x$ 将格点动量整数映射为物理相位。
      \end{block}
    \end{column}
    \begin{column}{0.45\textwidth}
      \begin{tabular}{@{}l@{}}
        \rowlabel{读入与涂抹伪代码} \\[0.08cm]
        $\text{Read raw f8 eigenvectors}$ \\[0.03cm]
        $\text{Infer }N_{\rm ev}\text{ from file size}$ \\[0.03cm]
        $\text{reshape}\to[N_{\rm ev},N_x^3,3]$ \\[0.03cm]
        $\text{complexify: }r+i\,q$ \\[0.03cm]
        $\text{truncate}\to[N_{\rm ev1},N_x^3,3]$ \\[0.03cm]
        $\text{phase}[x]\gets
        \exp[-2\pi i\,\vec p_{\rm smear}\!\cdot\!\vec x/N_x]$ \\[0.03cm]
        $\text{einsum: }v_{\rm mom}[v,x,a]=v[v,x,a]\,
        \text{phase}[x]$ \\[0.03cm]
        $\text{Output: GPU eigenvectors}$
      \end{tabular}
      \source{\code{2pt_Cg5gmu...gpu.py:152--186}；完整文件名见附录。}
    \end{column}
  \end{columns}
\end{frame}

% S09
\begin{frame}[t]{VVV 与 Wick 缩并：三夸克颜色反对称性落到两项五张量 contraction}
  \small
  \begin{center}
    \begin{tabularx}{0.95\textwidth}{@{}>{\raggedright\arraybackslash}p{0.18\linewidth}
      >{\raggedright\arraybackslash}X
      >{\raggedright\arraybackslash}p{0.27\linewidth}@{}}
      \toprule
      对象 & 代码动作 & 物理含义 \\
      \midrule
      VVV
        & 对三个颜色列做 $3!$ 个排列，符号为
          $(+,+,+,-,-,-)$，逐 $x$ 块累加
        & $\epsilon^{abc}$ 的颜色反对称质子插值结构，并带动量相位。 \\
      中间传播子
        & \code{CG5peram_u = contract(...)}
        & $\gamma_5$ 厄米关系/插值矩阵作用后的 perambulator。 \\
      直接项
        & \code{"abc,gjad,gjbe,ilcf,def->il"}
        & 三条传播线与源/汇 VVV 完成颜色、模、Dirac 缩并。 \\
      交换项
        & \code{"abc,glaf,gjbe,ijcd,def->il"}，前面减号
        & 费米子反对易与两条同味夸克交换的 Wick 拓扑。 \\
      输出
        & $C_{il}(t_s,t_i)$，再与 $\Gamma_\pm$ 缩并
        & 先得到 Dirac 矩阵，再选择正/负宇称通道。 \\
      \bottomrule
    \end{tabularx}
  \end{center}
  \source{\code{2pt_proton_Cg5gmu_L32x64_mom2_xdir_gpu.py:244--318,356--387}；
  重子插值物理式参照 \code{docs/pure_gluon_ope_distillation_20260815.tex:296--368}。}
\end{frame}

% S10
\begin{frame}[t]{DR $\gamma$ 矩阵、宇称投影与反周期边界必须保持同一约定}
  \small
  \begin{columns}[T,totalwidth=\textwidth]
    \begin{column}{0.49\textwidth}
      \begin{tabular}{@{}l@{}}
        \rowlabel{Dirac 约定} \\[0.08cm]
        $\gamma_0=\mathbf 1,\qquad
        \{\gamma_\mu,\gamma_\nu\}=2\delta_{\mu\nu}$ \\[0.04cm]
        $\gamma_5=\mathrm{diag}(1,1,-1,-1)$ \\[0.04cm]
        $\gamma_7=\gamma_3\gamma_1$ \\[0.04cm]
        \code{element="_Cg5g4"}:\\
        $\quad I_1= \gamma_7\gamma_4,\quad I_2=\gamma_7\gamma_4$ \\[0.04cm]
        $\text{插值作用：}\quad
        \tau\mapsto I_1\,\tau\,I_2$
      \end{tabular}
      \source{\code{gamma_matrix_cupy_DR.py:6--106}；
      \code{2pt_Cg5gmu...gpu.py:120--147}。}
    \end{column}
    \begin{column}{0.49\textwidth}
      \begin{tabular}{@{}l@{}}
        \rowlabel{宇称与边界算法} \\[0.08cm]
        $P_\pm=\frac12(1\pm\gamma_4)$ \\[0.04cm]
        $\text{Build }C_{il}(t_s,t_i)$ \\[0.04cm]
        $C_{+}\gets P_+^{\,li}C_{il},\qquad
          C_{-}\gets P_-^{\,li}C_{il}$ \\[0.04cm]
        $\text{if }t_s<t_i:\quad C_+(t_s,t_i)\gets-C_+(t_s,t_i)$ \\[0.04cm]
        $\text{if }t_s>t_i:\quad C_-(t_s,t_i)\gets-C_-(t_s,t_i)$ \\[0.04cm]
        $\text{then save }C_+\text{ (当前脚本)}$ \\[0.04cm]
        $\text{物理：三夸克反周期}$\\[-0.01cm]
        $\text{使 backward 态带反宇称结构}$
      \end{tabular}
      \source{\code{2pt_Cg5gmu...gpu.py:100--108,408--447}；
      反周期符号代码位于 \code{:414--425}。}
    \end{column}
  \end{columns}
  \vspace{0.12cm}
  \begin{block}{边界极限}
    若忽略绕回项，$t_s-t_i\ll N_t$ 时正宇称通道近似
    $C_2^{+}(t)\sim A_+e^{-E_+t}$；当时间间隔接近 $N_t$，不能把 backward
    反宇称贡献误当普通激发态。
  \end{block}
\end{frame}

% S11
\begin{frame}[t]{GPU/I/O：优化峰值显存而不改变物理收缩}
  \small
  \begin{center}
    \begin{tabularx}{0.95\textwidth}{@{}>{\raggedright\arraybackslash}p{0.19\linewidth}
      >{\raggedright\arraybackslash}p{0.34\linewidth}
      >{\raggedright\arraybackslash}X@{}}
      \toprule
      策略 & 实现 & 保持/改变的对象 \\
      \midrule
      后端
        & \code{cupy} 数组 + \code{opt_einsum.contract}
        & 保持 einsum 指标语义；把矩阵运算放到 GPU。 \\
      按源时间读入
        & 对每个 $t_{\rm source}$ 读取四个自旋文件
        & 峰值只承载当前源时间的 $\tau$，输出仍是完整
          $[N_t,N_t,4,4]$。 \\
      空间分块
        & VVV 的 \code{xi} 循环每次取 $N_x^2$ 个点
        & 结果逐块累加，避免一次生成过大的 $N_{\rm ev}^3N_x^3$ 中间量。 \\
      显存回收
        & \code{free_all_blocks()} + \code{synchronize()}
        & 释放缓存并确保异步 GPU 工作完成；不改变收缩公式。 \\
      原始 I/O
        & f8 实部/虚部交错；按文件大小反推 $N_{\rm ev}$
        & 兼容现有无元数据文件，但尺寸/端序错误会静默污染张量。 \\
      \bottomrule
    \end{tabularx}
  \end{center}
  \vspace{0.12cm}
  \begin{tabular}{@{}l@{}}
    $\text{raw bytes}\;\longrightarrow\;
    \text{reshape}\;\longrightarrow\;
    \text{complexify}\;\longrightarrow\;
    \text{truncate }N_{\rm ev1}\;\longrightarrow\;\text{GPU contraction}$ \\
    $\text{validation invariant: }\quad
    \text{element count}=
    2\times(\text{all declared dimensions})
    \quad\land\quad
    \text{complex dtype is explicit}$
  \end{tabular}
  \source{\code{input_output_4_cupy.py:17--66,105--199}；
  \code{2pt_gpu.py:250--307,391--392}。}
\end{frame}

% S12
\begin{frame}[t]{代码证据：驱动脚本把参数、方向、循环与落盘路径串起来}
  \small
  \lstset{basicstyle=\ttfamily\tiny}
  \begin{columns}[T,totalwidth=\textwidth]
    \begin{column}{0.48\textwidth}
      \begin{block}{OPE 驱动（直接引用）}
        \lstinputlisting[firstline=56,lastline=64]{../Calc_ope_unpol.py}
      \end{block}
      \source{\code{Calc_ope_unpol.py:56--64}；
      \code{link_dir} 映射到 Wilson 线方向。}
    \end{column}
    \begin{column}{0.48\textwidth}
      \begin{block}{2pt 关键输出（直接引用）}
        \lstinputlisting[firstline=414,lastline=424]{../2pt_proton_Cg5gmu_L32x64_mom2_xdir_gpu.py}
      \end{block}
      \source{\code{2pt_gpu.py:414--424}；
      先投影，再处理边界，再保存。}
    \end{column}
  \end{columns}
  \vspace{0.10cm}
  \begin{block}{解释}
    这两段分别证明：OPE 方向是显式输入并写入分方向目录；2pt 不是单个时间序列，
    而是先生成 $(t_s,t_i)$ 矩阵，再进行宇称投影与边界符号处理。代码片段是源文件直引，
    不是手抄伪代码。
  \end{block}
\end{frame}

% S13
\begin{frame}[t]{验证与风险：结构可读不等于数值链已闭合}
  \footnotesize
  \begin{center}
    \begin{tabularx}{0.95\textwidth}{@{}>{\raggedright\arraybackslash}p{0.21\linewidth}
      >{\raggedright\arraybackslash}p{0.27\linewidth}
      >{\raggedright\arraybackslash}X@{}}
      \toprule
      检查层级 & 当前证据 & 尚需执行的验证 \\
      \midrule
      语法/接口
        & 入口函数、张量 reshape、输出命名可静态追踪
        & 在目标环境执行 Python AST/语法检查与最小数组单元测试。 \\
      规范几何
        & Clover 四方向、$\epsilon$ 对偶、$U^\dagger$ 链路均有实现
        & 用单位规范场、单个非平凡 plaquette 检查 $F_{\mu\nu}=-F_{\nu\mu}$、
          $z=0$ 退化和方向翻转。 \\
      群论/Dirac
        & DR $\gamma$ 矩阵、$P_\pm$、两项 Wick 结构可追踪
        & 数值检查 $\{\gamma_\mu,\gamma_\nu\}$、$\gamma_5$ 厄米关系、
          projector 幂等性与正交性。 \\
      统计物理
        & 代码已分别保存 OPE 与 $C_2$
        & 真实组态上完成 ensemble covariance、真空扣除、$c_0$/
          矩阵元平台与归一化。 \\
      生产性能
        & 有分块、按源读入和 GPU 内存回收策略
        & 在真实设备记录峰值显存、耗时、MPI 负载均衡；本次没有硬件运行证据。 \\
      \bottomrule
    \end{tabularx}
  \end{center}
  \vspace{0.10cm}
  \begin{block}{最重要的风险排序}
    \rowlabel{R1} 无元数据原始二进制的 reshape/端序错误；
    \rowlabel{R2} Clover 实现未显式无迹投影；
    \rowlabel{R3} unpolarized 驱动中负方向数组的落盘被注释；
    \rowlabel{R4} 当前目录缺少真实数据，不能把静态审阅升级为数值正确性结论。
  \end{block}
  \source{\code{Operator.py:150--184}；
  \code{Calc_ope_unpol.py:116--141}；
  \code{input_output_4_cupy.py:44--66}。}
\end{frame}

% S14
\begin{frame}[t]{阶段结论：底层算法链可交接，但下游物理闭环仍需数据验证}
  \small
  \begin{center}
    \begin{tabularx}{0.95\textwidth}{@{}>{\raggedright\arraybackslash}p{0.20\linewidth}
      >{\raggedright\arraybackslash}p{0.37\linewidth}
      >{\raggedright\arraybackslash}X@{}}
      \toprule
      主张 & 证据 & 置信度与影响 \\
      \midrule
      OPE 算法结构成立
        & Clover $\to$ Wilson 线 $\to$ 颜色迹/空间和，
          正负方向与指标变体已编码
        & \verified\ 代码结构确证；可作为 OPE 计算模块交接。 \\
      2pt 算法结构成立
        & momentum-smearing distillation $\to$ VVV/peram
          $\to$ direct-exchange $\to P_\pm$
        & \verified\ 缩并路径确证；谱质量仍需实数拟合。 \\
      断连分解可用于组织流水线
        & OPE 与核子 2pt 使用独立输出文件/目录
        & \inferred\ 工程分解合理；组态统计组合不是乘法捷径。 \\
      最终胶子 PDF 已得到
        & 当前没有输入数据与下游 Fourier/匹配运行证据
        & \unverified\ 不作此结论；需补齐真实数据链。 \\
      \bottomrule
    \end{tabularx}
  \end{center}
  \vspace{0.12cm}
  \begin{tabular}{@{}l@{}}
    $\text{交接接口：}\quad
    \{O_{\mu\nu;\rho\sigma}(t,z),\;C_2^{\pm}(t_s,t_i,\vec P)\}
    \;\longrightarrow\;\text{ensemble analysis}$ \\
    $\text{验收门：}\quad
    \text{shape}\;\land\;\text{symmetry}\;\land\;\text{boundary}\;\land\;
    \text{covariance}\;\land\;\text{physical plateau}$
  \end{tabular}
  \source{综合 \code{Operator.py}、\code{Calc_ope_*.py}、
  \code{2pt_proton_Cg5gmu_*.py}；当前目录无生产数据输入。}
\end{frame}

% S15
\appendix
\begin{frame}[t]{附录 A：关键函数—物理对象映射表}
  \begin{center}
    \scriptsize
    \begin{tabularx}{0.95\textwidth}{@{}>{\raggedright\arraybackslash}p{0.31\linewidth}
      >{\raggedright\arraybackslash}p{0.29\linewidth}
      >{\raggedright\arraybackslash}X@{}}
      \toprule
      代码入口 & 张量/输出 & 物理角色 \\
      \midrule
      \code{plaquette_clover_new}
        & $[N_t,N_x,N_x,N_x,3,3]$
        & 局部 Clover 反厄米场强构造。 \\
      \code{plaquette_clover_all_new}
        & $[4,4,N_t,N_x,N_x,N_x,3,3]$
        & 遍历全部 Lorentz 指标，$\mu=\nu$ 置零。 \\
      \code{plaquette_clover_all_tilde}
        & 同形状
        & Levi--Civita 对偶张量。 \\
      \code{operators_new_z0_mu2}
        & $\mathbb C^{N_t}$
        & 任意 $(\mu_2,\nu_2)$ 的正方向非定域 OPE。 \\
      \code{operators_new_z0_mz_mu2}
        & $\mathbb C^{N_t}$
        & 负方向 Wilson 线变体。 \\
      \code{readin_eigvecs}
        & $[N_{\rm ev1},N_x^3,3]$
        & 空间 Laplacian 低模与动量涂抹。 \\
      \code{readin_peram}
        & $[N_t,4,4,N_{\rm ev1},N_{\rm ev1}]$
        & Dirac$\times$distillation 子空间传播子。 \\
      \code{VVV_Calc_cupy}
        & $[N_t,N_{\rm ev1},N_{\rm ev1},N_{\rm ev1}]$
        & 三夸克颜色反对称顶点。 \\
      \bottomrule
    \end{tabularx}
  \end{center}
  \source{\code{Operator.py:48--184,339--460}；
  \code{2pt_proton_Cg5gmu_L32x64_mom2_xdir_gpu.py:172--318}。}
\end{frame}

% S16
\begin{frame}[t]{附录 B：建议的最小数值验收顺序}
  \begin{center}
    \begin{tabular}{@{}l@{}}
      \rowlabel{Step 1 — 输入守卫} \quad
      构造小格点 raw f8，检查端序、实虚交错与 reshape 元素数 \\[0.06cm]
      \rowlabel{Step 2 — 规范几何} \quad
      $U_\mu=\mathbf 1\;\Longrightarrow\;F_{\mu\nu}=0,\;
      U^\dagger U=\mathbf 1\;\Longrightarrow\;$ Wilson 闭合稳定 \\[0.06cm]
      \rowlabel{Step 3 — Dirac/群论} \quad
      $\gamma_\mu^\dagger=\gamma_\mu,\;
      \{\gamma_\mu,\gamma_\nu\}=2\delta_{\mu\nu},\;
      P_\pm^2=P_\pm,\;P_+P_-=0$ \\[0.06cm]
      \rowlabel{Step 4 — OPE 边界} \quad
      $\delta z=0,\;\delta z=1,\;+\hat n,\;-\hat n$
      四个退化/方向用例逐项对照 \\[0.06cm]
      \rowlabel{Step 5 — 2pt 收缩} \quad
      小 $N_{\rm ev}$ 与显式循环结果对照 direct-exchange einsum \\[0.06cm]
      \rowlabel{Step 6 — 统计闭环} \quad
      \text{多组态上真空扣除、协方差、宇称通道和时间平台检查} \\[0.06cm]
      \rowlabel{Step 7 — 物理下游} \quad
      $\widetilde h(z,P_z)\to\widetilde g(x,P_z)\to g(x,\mu)$，记录匹配核与系统误差
    \end{tabular}
  \end{center}
  \vspace{0.14cm}
  \begin{block}{下一阶段可验收产物}
    一个可公开复现的单组态 smoke test、一个多组态协方差报告、一个
    Clover 无迹投影决策（修复或证明无需修复），以及 OPE$\times$2pt 到
    $\widetilde h$ 的明确归一化接口。
  \end{block}
  \source{验收顺序由本报告的物理不变量与代码接口推导；
  具体运行路径见 \code{donghx_code_analysis.tex}，算法段 119--1365。}
\end{frame}

\end{document}
